\section{Acids, Bases, Buffers, Titration, and Solubility Products}

\subsection{Ctrl+F keywords: pH, pOH, weak acid, \(K_a\), \(pK_a\), buffer, Henderson, phosphate, titration, \(K_{sp}\), common ion, oploeselighed}
Use for acid strength, pH calculations, buffers, phosphate systems, titration equivalence, solubility, precipitation, and common-ion tasks.

\subsection{Symbols and notation}
\referencetable{tab:acid:symbols}{Acid-base and solubility symbols}
\begin{tabular}{@{}>{$}l<{$} p{10.8cm}@{}}
\toprule
\text{Symbol} & \text{Meaning and usage}\\
\midrule
K_w,\ K_a,\ K_b & Water autoionization, acid dissociation, and base dissociation constants.\\
\mathrm{pH},\ \mathrm{pOH},\ pK_a,\ pK_b & Negative base-10 logarithms of hydronium, hydroxide, \(K_a\), and \(K_b\).\\
\ce{HA},\ \ce{A-},\ \ce{B},\ \ce{HB+} & Weak acid/conjugate base and weak base/conjugate acid notation.\\
c_0,\ x & Initial analytical concentration and amount dissociated per volume in an equilibrium setup.\\
K_{sp},\ Q_{sp},\ s & Solubility-product constant, current ion product, and molar solubility.\\
\mathrm M,\ \mathrm X & Generic cation/metal and anion in a sparingly soluble salt.\\
a,\ b,\ m,\ x & Stoichiometric subscripts or ionic charges in \(\mathrm M_a\mathrm X_b\); identified by context.\\
\bottomrule
\end{tabular}

\subsection{pH and acid-base constants}

\subsubsection{Water, pH, and pOH}
\begin{equation}
K_w=[\ce{H3O+}][\ce{OH-}]=1.0\times10^{-14}, \qquad
\mathrm{pH}=-\log[\ce{H3O+}], \qquad
\mathrm{pOH}=-\log[\ce{OH-}], \qquad
\mathrm{pH}+\mathrm{pOH}=14
\label{eq:acid:ph-poh-water}
\end{equation}
At \(298\,\mathrm K\), neutral water has pH \(7\). A decrease of one pH unit means ten times larger \([\ce{H3O+}]\).

\subsubsection{Strong acid and strong base pH}
\begin{equation}
[\ce{H3O+}]\approx \nu_{\ce{H+}}c_{\mathrm{acid}}, \qquad
[\ce{OH-}]\approx \nu_{\ce{OH-}}c_{\mathrm{base}}
\label{eq:acid:strong-acid-base}
\end{equation}
Strong acids and bases are treated as fully dissociated. Include stoichiometric coefficients; for example, \(\ce{Ba(OH)2}\) gives \(2\ce{OH-}\).

\subsubsection{Weak acid and weak base constants}
\begin{equation}
K_a=\frac{[\ce{A-}][\ce{H3O+}]}{[\ce{HA}]}, \qquad
K_b=\frac{[\ce{HB+}][\ce{OH-}]}{[\ce{B}]}
\label{eq:acid:weak-equilibrium-constants}
\end{equation}
\begin{equation}
K=10^{-pK}, \qquad pK=-\log K
\label{eq:acid:pk-conversion}
\end{equation}
Larger \(K_a\) or smaller \(pK_a\) means stronger acid. Larger \(K_b\) or smaller \(pK_b\) means stronger base.

\referencetable{tab:acid:common-strength-recognition}{Common acid and base identities for qualitative classification when no \(K_a\) data are supplied}
\begin{tabular}{@{}>{\raggedright\arraybackslash}p{3.0cm}>{\raggedright\arraybackslash}p{4.5cm}>{\raggedright\arraybackslash}p{4.5cm}@{}}
\toprule
Category & Course-level examples & Classification action\\
\midrule
Strong acids & \(\ce{HCl}\), \(\ce{HNO3}\), first proton of \(\ce{H2SO4}\) & Rank above listed weak acids in water.\\
Weak acids & \(\ce{HF}\), \(\ce{CH3COOH}\), \(\ce{HNO2}\), \(\ce{H3PO4}\) & Acidic, but weaker than the listed strong acids.\\
Bases / non-acids & \(\ce{NaOH}\), \(\ce{NH3}\), metal hydrides such as \(\ce{NaH}\) & Reject from an option requiring acids only.\\
Generally non-acid organic compounds & hydrocarbons and ordinary alcohols without an acidic group & Do not classify as common aqueous acids unless data state otherwise.\\
\bottomrule
\end{tabular}

\subsubsection{Conjugate acid-base relation}
\begin{equation}
K_aK_b=K_w, \qquad pK_a+pK_b=14
\label{eq:acid:conjugate-relation}
\end{equation}
Use this for conjugate acid/base pairs at \(298\,\mathrm K\).

\subsubsection{Weak acid and weak base approximation with exact fallback}
\begin{equation}
[\ce{H3O+}]\approx\sqrt{K_ac_0}, \qquad
[\ce{OH-}]\approx\sqrt{K_bc_0}
\label{eq:acid:weak-approximation}
\end{equation}
\begin{equation}
K=\frac{x^2}{c_0-x}
\quad\Longrightarrow\quad
x=\frac{-K+\sqrt{K^2+4Kc_0}}{2}
\label{eq:acid:weak-exact-quadratic}
\end{equation}
Use the approximation only when the computed \(x/c_0\leq 0.05\); otherwise use the exact positive root.

\subsection{Exam recipe: identify acids and rank acid strength}
\textbf{Use when:} answer options list chemical formulas and ask which is an acid, which set contains only acids or weak acids, or which acid is strongest without supplying numerical constants.

\textbf{Given / Find:} Given formulas and possibly \(K_a\) or \(pK_a\) data; find acid/non-acid classification or relative acid strength.

\textbf{Required references:} Table~\ref{tab:acid:common-strength-recognition} and Equation~\eqref{eq:acid:pk-conversion} when numerical \(K_a\) or \(pK_a\) data are given.

\textbf{Steps:}
\begin{enumerate}
    \item Match each listed species to Table~\ref{tab:acid:common-strength-recognition}; reject an ``all acids'' or ``all weak acids'' option as soon as it contains a listed base or non-acid.
    \item For qualitative ranking without data, rank a listed strong acid above a listed weak acid; compare weak acids further only when course data or a supplied table supports it.
    \item When \(K_a\) values are supplied, rank acids by decreasing \(K_a\); when \(pK_a\) values are supplied, rank them by increasing \(pK_a\), using Equation~\eqref{eq:acid:pk-conversion}.
    \item State the classification or selected option and identify the specific species or constant comparison that decides it.
\end{enumerate}

\textbf{Checks:} Strong acid means essentially complete dissociation for ordinary introductory calculations; smaller \(pK_a\) must correspond to larger \(K_a\).

\textbf{Common traps:} assuming every compound containing hydrogen is an acid; classifying \(\ce{NH3}\) as an acid because it contains hydrogen; ranking weak acids without supplied or established comparative data.

\subsection{Exam recipe: weak acid or weak base pH with approximation and exact fallback}
\textbf{Use when:} a weak acid/base concentration and \(K_a\), \(K_b\), \(pK_a\), or \(pK_b\) are given without a pre-existing buffer pair.

\textbf{Given / Find:} given \(c_0\) and an acid/base constant; find \([\ce{H3O+}]\) or \([\ce{OH-}]\), then pH.

\textbf{Required references:} equations~\eqref{eq:acid:ph-poh-water}, \eqref{eq:acid:weak-equilibrium-constants}, \eqref{eq:acid:pk-conversion}, \eqref{eq:acid:conjugate-relation}, \eqref{eq:acid:weak-approximation}, and \eqref{eq:acid:weak-exact-quadratic}.

\textbf{Steps:}
\begin{enumerate}
    \item Select the equilibrium constant by using \(K_a\) for \(\ce{HA}\) or \(K_b\) for \(\ce{B}\), converting a supplied \(pK\) with Equation~\eqref{eq:acid:pk-conversion}, and using Equation~\eqref{eq:acid:conjugate-relation} only when the conjugate constant is supplied.
    \item Estimate dissociation by calculating \(x_{\mathrm{approx}}=\sqrt{Kc_0}\) from equation~\eqref{eq:acid:weak-approximation}, interpreting \(x\) as \([\ce{H3O+}]\) for an acid and \([\ce{OH-}]\) for a base.
    \item Decide whether the estimate is acceptable by calculating \(x_{\mathrm{approx}}/c_0\); retain it when the ratio is at most \(0.05\), otherwise calculate \(x\) from equation~\eqref{eq:acid:weak-exact-quadratic}.
    \item Report pH by evaluating \(-\log[\ce{H3O+}]\) for an acid or by evaluating pOH from \([\ce{OH-}]\) and subtracting it from \(14\) for a base using equation~\eqref{eq:acid:ph-poh-water}.
\end{enumerate}

\textbf{Checks:} \(0<x<c_0\); a weak acid solution should be acidic and a weak base solution basic under ordinary concentrations; the approximation and exact result should be close when the 5\% test passes.

\textbf{Common traps:} using the square-root approximation after it fails the 5\% test; reporting pOH as pH for a weak base; mixing \(K_a\) with a base calculation without converting through \(K_w\).

\subsection{Buffers and titration}

\subsubsection{Henderson-Hasselbalch buffer equation}
\begin{equation}
\mathrm{pH}=pK_a+\log\left(\frac{[\ce{A-}]}{[\ce{HA}]}\right)
=pK_a+\log\left(\frac{n_{\ce{A-}}}{n_{\ce{HA}}}\right)
\label{eq:acid:henderson-hasselbalch}
\end{equation}
The mole ratio form applies when acid and base occupy the same final solution volume.

\subsubsection{Buffer after adding strong acid or base}
\begin{equation}
\ce{A- + H+ -> HA}, \qquad
\ce{HA + OH- -> A- + H2O}
\label{eq:acid:buffer-neutralization}
\end{equation}
Complete the stoichiometric neutralization before applying the buffer equation.

\subsubsection{Phosphoric acid phosphate system}
\begin{equation}
\begin{aligned}
\ce{H3PO4 &<=> H+ + H2PO4-} && pK_{a1}\approx2.1,\\
\ce{H2PO4- &<=> H+ + HPO4^2-} && pK_{a2}\approx7.2,\\
\ce{HPO4^2- &<=> H+ + PO4^3-} && pK_{a3}\approx12.4
\end{aligned}
\label{eq:acid:phosphate-pka}
\end{equation}
Choose the conjugate pair whose \(pK_a\) is closest to the target or observed pH.

\subsubsection{Titration equivalence point}
\begin{equation}
n_{\mathrm{acid\ equivalents}}=n_{\mathrm{base\ equivalents}}
\label{eq:acid:titration-equivalence}
\end{equation}
At equivalence, the pH depends on the acid-base properties of the salt remaining after neutralization.

\subsubsection{Complete neutralization of a polyprotic acid}
\begin{equation}
\ce{H3PO4 + 3OH- -> PO4^3- + 3H2O}, \qquad
n_{\ce{OH-}}=3n_{\ce{H3PO4}}
\label{eq:acid:phosphoric-complete-neutralization}
\end{equation}
To form fully deprotonated phosphate, supply one mole of hydroxide for each acidic proton removed.

\subsection{Exam recipe: acid-base stoichiometry and titration amount}
\textbf{Use when:} the task asks for the amount or volume of strong acid/base needed to neutralize an acid/base or to form a specified phosphate species.

\textbf{Given / Find:} Given concentrations, volumes, and target neutralization product; find moles or volume of titrant.

\textbf{Required references:} Equations~\eqref{eq:acid:titration-equivalence} and \eqref{eq:acid:phosphoric-complete-neutralization}; use \(n=cV\) from Equation~\eqref{eq:aqueous:concentration-dilution}.

\textbf{Steps:}
\begin{enumerate}
    \item Identify how many protons or hydroxides must react per formula unit to produce the named target species; for sodium phosphate from \(\ce{H3PO4}\), use the 3:1 hydroxide-to-acid ratio in Equation~\eqref{eq:acid:phosphoric-complete-neutralization}.
    \item Calculate initial acid or base moles from its concentration and volume.
    \item Multiply by the required stoichiometric equivalents to find titrant moles, using Equation~\eqref{eq:acid:titration-equivalence}.
    \item Divide required titrant moles by titrant concentration to obtain the volume.
\end{enumerate}

\textbf{Checks:} Conversion to fully deprotonated \(\ce{PO4^3-}\) requires three hydroxide equivalents per \(\ce{H3PO4}\); use liters in \(n=cV\).

\textbf{Common traps:} stopping at the first equivalence point when the target salt requires full deprotonation; using milliliters without converting to liters.

\subsection{Exam recipe: buffer pH}
\textbf{Use when:} both a weak acid and its conjugate base are present, including a phosphate buffer or a buffer after addition of strong acid/base.

\textbf{Given / Find:} given pair amounts, \(pK_a\), and optionally added acid/base; find buffer pH or the required ratio.

\textbf{Required references:} equations~\eqref{eq:acid:henderson-hasselbalch}, \eqref{eq:acid:buffer-neutralization}, and \eqref{eq:acid:phosphate-pka} for phosphate buffers.

\textbf{Steps:}
\begin{enumerate}
    \item Identify the conjugate pair by matching species that differ by one proton, and for phosphate choose the pair in equation~\eqref{eq:acid:phosphate-pka} with \(pK_a\) nearest the relevant pH.
    \item Convert the supplied amounts to moles by multiplying each concentration by its solution volume before any mixing or reaction.
    \item Account for added strong acid or base by subtracting and adding moles according to equation~\eqref{eq:acid:buffer-neutralization} until the limiting strong reagent is consumed.
    \item Calculate pH by inserting the remaining mole ratio \(n_{\ce{A-}}/n_{\ce{HA}}\) into equation~\eqref{eq:acid:henderson-hasselbalch}, or rearrange that equation to find a requested ratio.
\end{enumerate}

\textbf{Checks:} equal conjugate-pair amounts give \(\mathrm{pH}=pK_a\); a valid buffer retains positive amounts of both pair members; adding acid must lower the pH.

\textbf{Common traps:} applying Henderson-Hasselbalch before neutralizing added strong reagent; selecting the wrong phosphate \(pK_a\); using concentrations from unequal pre-mixing volumes instead of post-reaction mole ratios.

\subsection{Solubility products}

\subsubsection{Solubility product expression}
\begin{equation}
\mathrm{M}_a\mathrm{X}_b(s)\rightleftharpoons a\,\mathrm{M}^{m+}+b\,\mathrm{X}^{x-}
\label{eq:acid:dissolution-reaction}
\end{equation}
\begin{equation}
K_{sp}=[\mathrm{M}^{m+}]^a[\mathrm{X}^{x-}]^b
\label{eq:acid:ksp-expression}
\end{equation}
The pure solid is omitted. Powers come from dissolution stoichiometry.

\subsubsection{Molar solubility patterns}
\referencetable{tab:acid:solubility-patterns}{Common molar-solubility substitutions}
\begin{tabular}{@{}l l l@{}}
\toprule
\text{Solid} & \text{Ion concentrations from molar solubility \(s\)} & \text{\(K_{sp}\) relation}\\
\midrule
\(\mathrm{MX}\) & \([\mathrm M^+]=s,\ [\mathrm X^-]=s\) & \(K_{sp}=s^2\)\\
\(\mathrm{MX_2}\) & \([\mathrm M^{2+}]=s,\ [\mathrm X^-]=2s\) & \(K_{sp}=4s^3\)\\
\(\mathrm{M(OH)_2}\) & \([\mathrm M^{2+}]=s,\ [\ce{OH-}]=2s\) & \(K_{sp}=4s^3\)\\
\bottomrule
\end{tabular}

\subsubsection{Ion product and precipitation}
\begin{equation}
Q_{sp}=[\mathrm{M}^{m+}]^a[\mathrm{X}^{x-}]^b
\label{eq:acid:qsp-expression}
\end{equation}
\referencetable{tab:acid:precipitation-test}{Precipitation decision from \(Q_{sp}\) and \(K_{sp}\)}
\begin{tabular}{@{}c l@{}}
\toprule
\text{Comparison} & \text{State or direction}\\
\midrule
\(Q_{sp}<K_{sp}\) & Unsaturated; additional solid can dissolve\\
\(Q_{sp}=K_{sp}\) & Saturated solution at equilibrium\\
\(Q_{sp}>K_{sp}\) & Supersaturated; precipitation is favored\\
\bottomrule
\end{tabular}

\subsubsection{Common-ion solubility}
\begin{equation}
K_{sp}=[\mathrm{M}^{m+}][\mathrm{X}^{x-}], \qquad
[\mathrm{X}^{x-}]_{\mathrm{initial}}\gg s
\quad\Longrightarrow\quad
s\approx\frac{K_{sp}}{[\mathrm{X}^{x-}]_{\mathrm{initial}}}
\label{eq:acid:common-ion-approximation}
\end{equation}
A common ion lowers solubility by shifting dissolution toward the solid.

\subsection{Exam recipe: solubility product, common ion, and hydroxide pH}
\textbf{Use when:} the task asks for molar solubility, precipitation, saturated-solution pH, or common-ion effects.

\textbf{Given / Find:} given a salt formula, \(K_{sp}\), and optionally initial ion concentrations; find \(s\), \(Q_{sp}\), precipitation direction, or pH for a hydroxide.

\textbf{Required references:} equations~\eqref{eq:acid:dissolution-reaction}, \eqref{eq:acid:ksp-expression}, \eqref{eq:acid:qsp-expression}, \eqref{eq:acid:common-ion-approximation}, and \eqref{eq:acid:ph-poh-water}; Tables~\ref{tab:acid:solubility-patterns} and \ref{tab:acid:precipitation-test}.

\textbf{Steps:}
\begin{enumerate}
    \item Write the dissolution equation by splitting one formula unit into ions and preserving each stoichiometric subscript as an ion coefficient.
    \item Build the \(K_{sp}\) or \(Q_{sp}\) expression by raising each dissolved ion concentration to its coefficient as in equation~\eqref{eq:acid:ksp-expression}.
    \item Calculate molar solubility by replacing each newly formed ion concentration with its coefficient times \(s\), using Table~\ref{tab:acid:solubility-patterns} directly for the listed patterns and adding any given initial common-ion concentration before solving.
    \item Decide whether precipitation occurs by calculating \(Q_{sp}\) from current concentrations and comparing it with \(K_{sp}\) using Table~\ref{tab:acid:precipitation-test}.
    \item Find hydroxide-solution pH by converting the solved \(s\) into \([\ce{OH-}]\) through the dissolution coefficient, then calculating pOH and pH with equation~\eqref{eq:acid:ph-poh-water}.
\end{enumerate}

\textbf{Checks:} a common ion must reduce \(s\) relative to pure water; for \(\mathrm{M(OH)_2}\), \([\ce{OH-}]=2s\); precipitation begins at \(Q_{sp}=K_{sp}\).

\textbf{Common traps:} including the solid in \(K_{sp}\); missing powers generated by stoichiometric coefficients; setting \([\ce{OH-}]=s\) for a dihydroxide; applying the common-ion approximation when the initial ion concentration is not much greater than \(s\).
